\begin{abstract}
Modular foundation servers offer enterprises an avenue to reduce costs associated with large language models. By keeping a single, fixed foundation model in memory, LoRA matrices can be swapped in to adapt the model to new domains. As such, compute costs decouple from domain adaptation costs. However, this introduces security concerns: inference involves running a common base model with an ever-increasing number of adapters, many of which originate from open source repositories. As such, each adapter in the inference pipeline poses a potential attack vector.

In this paper we consider two vulnerabilities. First, base model auditing (RQ1). Due to variations in activation scaling, normalisation, and attention mechanisms between architectures, detectors cannot distinguish between benign distribution shifts and adversarial poisoning in an uncalibrated manner. Thus, many attacks go undetected or many false positives occur. Second, adapter loading (RQ2). Adapters downloaded from public repositories typically undergo no static verification before entering inference on GPUs, exposing inference to unsafe adapters, stealthy low-rank trojan attacks and crashing loaders.

We present CALB-Shield. To address RQ1, we propose Parent-Anchored Behavioural Normalisation. We gather outputs for each of 30 probes to form a fingerprint with dimensionality 180. We achieve numerical stability in our fingerprints (coefficient of variation: 0.0446\% across 150 runs on hardware). For each potential checkpoint, we compare to the declared parent model of the checkpoint rather than a global baseline. To remove any bias from model transfers between architectures, we achieve perfect detection of physical backdoors (100.0\%) for 5 different classifier types. We achieve perfect macro-F1 scores (1.0) using 5-fold Leave-One-Pretrained-Out cross-validation, and detect no false positives in benign parameter drift.

To address RQ2, we find a trade-off between inference safety and utility. Rigorous checks during inference can cause performance bottlenecks but leaving checks out causes exposure. We perform rigorous checks in a pre-inference pipeline on the CPU. Checks include cryptographically secure AIBOM verification, rank-aware spectral analysis via QR-SVD, trigger inversion, and safety checks. Spectral analysis enables decomposition of low-rank matrices in less than 0.25 seconds (5000$\times$ faster than dense SVD). To maintain utility in open source communities, we do not outright reject adapters exhibiting ambiguous fine-tuning behaviour. Instead we admit them into a quarantined environment under observation, and pair them with an asynchronous streaming agent during inference.

On a benchmark with 71 targets spanning 4 families of models, we achieve perfect detection of structural incompatibility and malicious adapters. (28/28 malicious adapters detected with 0 token leakage). We admitted 96.88\% of clean community adapters (98.44\% balanced accuracy; Zou-Donner 95\% CI [89.71\%, 99.72\%]). The sidecar for base model auditing has an average latency of 0.0412ms (41.2\,$\mu$s) (0.008552\% overhead in relation to typical GPU inference).

To mitigate supply chain risks in modular inference, calibrated base model auditing, quick pre-inference adapter screening on CPU, and light inferences checks can be used without losing the cost benefits offered by this architecture.
\end{abstract}

\begin{IEEEkeywords}
Large Language Models, Backdoor Detection, Cross-Architecture Transferability, PEFT / LoRA Security, Supply Chain, Runtime Guardrails, AIBOM.
\end{IEEEkeywords}
